\documentclass[a4paper, 12pt]{article}

%-----------------------------------------------------
%                      PACKAGES
% ----------------------------------------------------
\usepackage[T1]{fontenc}
\usepackage[latin9]{inputenc}
\usepackage{fancyhdr}
\usepackage{geometry}
\usepackage[colorlinks=true, linkcolor=blue,citecolor=blue, urlcolor=blue]{hyperref}
\usepackage{indentfirst}
\usepackage{graphicx}
\usepackage{float}
\usepackage{setspace}
\usepackage{amsmath}
\usepackage{amssymb}
\usepackage{multirow}
\usepackage[table,xcdraw]{xcolor}
\usepackage{colortbl}
\usepackage{tabularx,booktabs}
\usepackage{enumerate}
\usepackage{titlesec}
\usepackage{caption}
\usepackage{subcaption}
\usepackage{nicefrac, xfrac}
\usepackage{subfiles}
\usepackage{mathtools}
\usepackage{bm}
\usepackage{tikz}
\usepackage[mode=buildnew]{standalone}
\usepackage{systeme}
\usepackage{stmaryrd}
\usepackage{import}
\usepackage[english]{babel} 
\usepackage{times}
\usepackage{pgfplots}
\usepackage{makecell}
\usepackage{threeparttable}
\usepackage{ragged2e}
\usepackage{tocloft}
\usepackage{booktabs}
\usepackage[acronym]{glossaries}
\usepackage{pdfpages}
\usepackage{lipsum}
\usepackage{hyphenat} 
\usepackage{xspace}
\usepackage{abntex2cite}
\usepackage{pdflscape}
\usepackage{pythonhighlight}
\usepackage{listings}
\usepackage{makecell}
\usepackage{longtable}
\usepackage{adjustbox}

%-----------------------------------------------------
%               DOCUMENT SETTINGS
% ----------------------------------------------------
\geometry{a4paper,top=30mm,bottom=20mm,left=30mm,right=20mm}

\pgfplotsset{compat=1.17} 

\numberwithin{equation}{section}
\counterwithin{figure}{section}
\counterwithin{table}{section}
% \numberwithin{lstlisting}{section}

% Defining spaces between lines
\setstretch{1.5} 

\graphicspath{{Figures/}}

\tikzset{
	font={\fontsize{11pt}{12}\selectfont}}

\setcounter{secnumdepth}{4}
\setcounter{tocdepth}{4}

% defining the height of table cells
\renewcommand{\arraystretch}{1.4}

% defining the gradient, divergence, and curl operators
\newcommand{\grad}[1]{\nabla #1}
\renewcommand{\div}[1]{\nabla \cdot #1}
\newcommand{\curl}[1]{\nabla \times #1}

\def\textcite#1{\citeauthoronline{#1} \cite{#1}}
%-----------------------------------------------------
%              PRE TEXTUAL ELEMENTS
% ----------------------------------------------------
\newenvironment{epigrafe}{\newpage\mbox{}\vfill\hfill\begin{minipage}[t]{0.5\textwidth}}
{\end{minipage}\newpage}

%-----------------------------------------------------
%              PYTHON CODE SETTINGS
% ----------------------------------------------------
\definecolor{codegreen}{rgb}{0,0.6,0}
\definecolor{codegray}{rgb}{0.5,0.5,0.5}
\definecolor{codepurple}{rgb}{0.58,0,0.82}
\definecolor{backcolour}{rgb}{0.95,0.95,0.92}

\lstdefinestyle{mystyle}{
    backgroundcolor=\color{backcolour},   
    commentstyle=\color{codegreen},
    keywordstyle=\color{magenta},
    numberstyle=\tiny\color{codegray},
    stringstyle=\color{codepurple},
    basicstyle=\ttfamily\footnotesize,
    breakatwhitespace=false,         
    breaklines=true,                 
    captionpos=b,                    
    keepspaces=true,                 
    numbers=left,                    
    numbersep=5pt,                  
    showspaces=false,                
    showstringspaces=false,
    showtabs=false,                  
    tabsize=2
}

\lstset{style=mystyle}
%-----------------------------------------------------
%                     DOCUMENT 
% ----------------------------------------------------
\begin{document}
% ------- WORK INFORMATION -------
\author{Carlos Henrique Chama Puga }
\newcommand{\RA}{195416}

\title{List 5 \\ Numerical Solutions of Nonlinear Equations}
\newcommand{\theme}{Numerical Solutions of Nonlinear Equations}

\newcommand{\Uni}{{Universidade Estadual de Campinas}}
\newcommand{\Fac}{{Faculdade de Engenharia Civil, Arquitetura e Urbanismo}}

\newcommand{\advisor}{Porf. Dr. Philippe Devloo}
\newcommand{\coadvisor}{Dr. Giovane Avancini}

\newcommand*{\workyear}{2024}

\makeatletter

% ------- PRE TEXTUAL PAGES -------
% Could not find another way to remove the page number from the first pages
\fancypagestyle{plain}{
    \fancyhf{}% Limpa todos os campos
    \fancyfoot[C]{}%
    \renewcommand{\headrulewidth}{0pt}%
}

\pagestyle{plain}
\def\logos{
    \noindent
    \raisebox{-.5\height}{\includegraphics[width=2.2cm]{Figures/logo-unicamp.pdf}}

    \vspace*{1.5cm}
    
    \noindent
    \begin{center} \large
        \MakeUppercase{\Uni}\\
        \Fac\\
    \end{center}
}

\def\openningpage{
  \logos
  \vskip 35mm
  \begin{center}
    \Large
      {\bf \@title}
  \end{center}
  \vskip 25mm
  \begin{flushright}
    \large
    {\textbf{Aluno:} \\ \@author - \RA}
    \vskip 10mm
    {{\bf{Docentes}}: \\
     \advisor \\
     \coadvisor}
  \end{flushright}
    \vfill
    \large
  \begin{center}
    Campinas\\\workyear
  \end{center}
}

\openningpage % Cover page

% Table of contents
\newpage
\tableofcontents 

% From here on, the page number is shown
% Chapter first page settings
\newpage
\fancypagestyle{plain}{
    \fancyhf{}% Limpa todos os campos
    \fancyhead[R]{\thepage}%
    \renewcommand{\headrulewidth}{0pt}%
}

\fancypagestyle{headings}{%
    \fancyhf{}% Limpa todos os campos
    \fancyhead[L]{\theme}% Nome do trabalho à esquerda
    \fancyhead[R]{\thepage}% Numero da página à direita
    \renewcommand{\headrulewidth}{1pt}%
}

% ------- CHAPTERS -------
\pagestyle{headings}
\section{Introduction} \label{sec:introduction}
It is possible to think of a matrix as a linear transformation applied to a vector. When the matrix is applied to a vector and the resulting vector is parallel to the original vector, this vector is said to be an eigenvector of the matrix so that 
\begin{equation}
    A\mathbf{v} = \lambda \mathbf{v},
    \label{eq:eigenvalue}
\end{equation}
in which $A$ is the matrix, $\mathbf{v}$ is the eigenvector, and $\lambda$ is the eigenvalue. 

Equation \eqref{eq:eigenvalue} can be rewritten as
\begin{equation}
    (A - \lambda I)\mathbf{v} = 0,
    \label{eq:eigenvalue_analytical}
\end{equation}
which means that, for any nonzero vector $\mathbf{v}$, Eq. \eqref{eq:eigenvalue_analytical} holds only if the determinant of the matrix $A - \lambda I$ is zero. In other words, matrix $A - \lambda I$ is singular. This is the characteristic equation of the matrix $A$ and can be used to find the eigenvalues of the matrix. 

Naturally, the eigenvalues of a matrix are the roots of the characteristic equation. However, a generalization of the characteristic equation is not always possible, and therefore, its analytical solution is not always feasible. 

The Power Method is an iterative method that can be used to find the most prominent eigenvalue of a matrix. The method is simple to implement and can be used to find both the eigenvalue and the eigenvector of the matrix. 

For a given vector $\mathbf{y}$ linearly independent of $\mathbf{v}$, it is possible to write
\begin{equation}
    \mathbf{y} = \sum_{i=1}^{n}\beta_i\mathbf{v}_i,
    \label{eq:linear_combination}
\end{equation}
in which $\mathbf{v}_i$ are the eigenvectors of the matrix $A$. 

Multiplying Eq. \eqref{eq:linear_combination} by $A^k$ yields
\begin{equation}
    A^k\mathbf{y} = \sum_{i=1}^{n}\beta_i\lambda_i^k\mathbf{v}_i,
    \label{eq:eigenvalue_power}
\end{equation}
and factoring Eq. \eqref{eq:eigenvalue_power} by $\lambda_1^k$ gives
\begin{equation}
    {A^k\mathbf{y}} = {\lambda_1^k}\sum_{i=1}^{n}\beta_i\left(\frac{\lambda_i}{\lambda_1}\right)^k\mathbf{v}_i.
    \label{eq:eigenvalue_power_factor}
\end{equation}

Since $\lambda_1$ is the most prominent eigenvalue, the term $\left(\frac{\lambda_i}{\lambda_1}\right)^k$ tends to zero for $i \neq 1$. Therefore, the term $\sum_{i=1}^{n}\beta_i\left(\frac{\lambda_i}{\lambda_1}\right)^k\mathbf{v}_i$ tends to $\beta_1\lambda_1^k\mathbf{v}_1$ as $k$ tends to infinity. 

This is the main idea behind the Power Method. The method starts with an initial guess $\mathbf{v}_0$. The matrix $A$ is multiplied by $\mathbf{v}_0$ to find the next eigenvector
\begin{equation}
    \mathbf{y} = A\mathbf{v}_{i-1},
    \label{eq:y}
\end{equation}
the norm of the result is calculated and considered the new eigenvalue
\begin{equation}
    \lambda_i = \left\|\mathbf{y}\right\|,
\end{equation}
the new eigenvector is the previous result divided by the new eigenvalue
\begin{equation}
    \mathbf{v}_i = \frac{\mathbf{y}}{\lambda_i},
\end{equation}

This process is repeated until the error between the current and previous eigenvalues is smaller than a given precision. Since Eq. \eqref{eq:y} is performed at each iteration, the limit when $k$ tends to infinity is approached.

The matrix $A$ can also be decomposed into the matrices $\Lambda$ and $Q$ so that
\begin{equation}
    A = Q\Lambda Q^{-1},
    \label{eq:decomposition}
\end{equation}
where $\Lambda$ is a diagonal matrix with the eigenvalues of $A$ 
\begin{equation*}
    \Lambda = \begin{bmatrix}
        \lambda_1 & 0 & \cdots & 0 \\
        0 & \lambda_2 & \cdots & 0 \\
        \vdots & \vdots & \ddots & \vdots \\
        0 & 0 & \cdots & \lambda_n
    \end{bmatrix},
\end{equation*}
and $Q$ is a matrix with the eigenvectors of $A$
\begin{equation*}
    Q = \begin{bmatrix}
        \mathbf{v}_1 & \mathbf{v}_2 & \cdots & \mathbf{v}_n
    \end{bmatrix}.
\end{equation*}

In the next sections, the power method is implemented and the results are discussed. The following bibliography is referred to in this work: \cite{de2000metodos,burden1997numerical}.
\section{Code Implementation}\label{sec:code_implementation}
This section presents the implementation of the classes used to solve the exercises proposed. For solving IVP, the RungeKutta class is developed. For BVP, the Galerkin class is implemented. Finally, for the Least Square Method, the LeastSquare class is created.

Code \ref{code:rk} shows the main function for the Runge-Kutta problem. 
\begin{lstlisting}[caption={Main function for the Runge-Kutta problem},label={code:rk},language=python]
def main_rk()->None:
    ft = lambda t, y: 1 + (t - y) ** 2
        
    rk = RungeKutta(ft, 2, 3, 0.1, 1)
    
    # 3/8 rule 4th order Runge-Kutta method
    rk.SetButcherTableau(method = "ThreeEights")
    rk.Run()

    print(rk.sol)

    rk.WriteResults("List7/Results_RK.txt")

    return
\end{lstlisting}

On line 2, the right-hand side of the differential equation is defined. On line 4, the RungeKutta object is created with the function, initial and final times, time step, and initial condition. The Butcher Tableau is set on line 7, and the Run method is called on line 8. The results are printed on line 10, and written to a file on line 12.

The main function for the Least Squares Method is shown in Code \ref{code:ls}
\begin{lstlisting}[caption={Main function for the Least Squares Method},label={code:ls},language=python]
def LeastSquareMethod()->None:
    data = [xi, yi]

    squares = LeastSquare(*zip(*data), "NonLinear")
    squares.Run()

    print(f"{squares.alpha=}")
    print(f"{squares.approx_solution=}")
    print(f"{squares.errors=}")
    print(f"{squares.total_error=}")

    return
\end{lstlisting}

Where line 2 defines the set of data points. The LeastSquare object is created on line 4, and the run method is called on line 5. The results are printed on lines 7 to 10.


\subsection{The RungeKutta Class}\label{subsec:rungekutta_class}
The RungeKutta class' attributes are displayed in Code \ref{code:rungekutta_attributes}. 
\begin{lstlisting}[caption={Attributes of the RungeKutta class},label={code:rungekutta_attributes},language=python]
@dataclass
class RungeKutta:
    ft: callable
    t0: float
    tf: float
    Dt: float
    y0: float

    ButcherTableau: list[float] = field(init=False, default_factory=list)
    c: list[float] = field(init=False, default_factory=list)
    a: list[list[float]] = field(init=False, default_factory=list)
    b: list[float] = field(init=False, default_factory=list)

    sol: list[float] = field(init=False, default_factory=list)
    step: float = field(init=False, default=0)
\end{lstlisting}

Where f\_t is the function to be solved, t0 and tf are the initial and final times, Dt is the time step, and y0 is the initial condition. The Butcher Tableau is a list of lists containing the coefficients of the Runge-Kutta method. The c, a, and b lists are the coefficients of the Butcher Tableau. The sol list stores the solution to the problem, and the step attribute is used to store the current time step.

As methods, the class has the SetButcherTableau, ConstantK, and Run. Other methods such as the EulerMethod, RK2, RK4, ThreeEightsRule, and WriteResults are also implemented. 

The SetButcherTableau method is shown in Code \ref{code:setbutchertableau_method}
\begin{lstlisting}[caption={SetButcherTableau method},label={code:setbutchertableau_method},language=python]
def SetButcherTableau(self, **var)->None:
    butcher = {"euler": self.EulerMethod(), "rk2": self.RK2(), "rk4": self.RK4(), "ThreeEights": self.ThreeEighthRule()}

    if 'method' in var:
        butcher[var['method']]

    elif 'ButcherTableau' in var:
        self.ButcherTableau = var['ButcherTableau']

    else: 
        raise ValueError("Invalid Butcher Tableau")
\end{lstlisting}

This method sets the Butcher Tableau according to the method chosen. The Euler method, Runge-Kutta 2nd order, Runge-Kutta 4th order, and Three-Eighths Rule are implemented. The user can also set the Butcher Tableau manually, by passing the desired rule. 

The ConstantK method is shown in Code \ref{code:constantk_method}
\begin{lstlisting}[caption={ConstantK method},label={code:constantk_method},language=python]
def ConstantK(self, index, t, y, k)->float:
    a = t + self.c[index] * self.Dt
    b = y + self.Dt * sum([self.a[index][j] * k[j] for j in range(index)])

    return self.ft(a, b)
\end{lstlisting}

This method is responsible for evaluating the intermediate values of the Runge-Kutta method. The index is the stage of the method, t is the current step time, y is the current value of the function, and k is the list of intermediate values. The return is the function evaluated at the intermediate point.

Finally, the Run method is shown in Code \ref{code:run_method}
\begin{lstlisting}[caption={Run method},label={code:run_method},language=python]
def Run(self)->None:
    t = self.t0
    y = self.y0

    self.c, self.a, self.b = self.ButcherTableau

    k = []
    for _ in range(self.step):
        for i in range(len(self.c)):
            k.append(self.ConstantK(i, t, y, k))

        y += self.Dt * sum([self.b[j] * k[j] for j in range(len(k))])
        t += self.Dt

        k.clear()
        self.sol.append((t, y))
\end{lstlisting}

The Run method performs the Runge-Kutta method for a given Butcher Tableau. On lines 2 and 3, the initial values of t and y are set. The Butcher Tableau coefficients (c, a, and b) are unpacked on line 5. 

From line 8 on, the method iterates over the number of steps. For each step, the intermediate values are calculated on line 10. The new value of y is calculated on line 12, and the new value of t is calculated on line 13. The intermediate values are cleared on line 15, and the t and y values are appended to the sol list on line 16.

\subsection{The Galerkin Class}\label{subsec:galerkin_class}

\subsection{The LeastSquare Class}\label{subsec:leastsquares_class}
\section{Results} \label{sec:results}
\section{Conclusions}\label{sec:conclusions}
The present work aimed to solve the ordinary differential equation using the Runge-Kutta method. An algorithm is developed to read a given Butcher Tableau and approximate the solution for the ODE. The main difference between initial value problems and boundary value problems is discussed, and for the latter, a Galerkin algorithm is also implemented.

For the Galerkin method, a few comments are in order. Since the method is based on the integration of functions, the number of integration points when using the Gaussian Quadrature is key to the accuracy of the results. During the elaboration of the code, it was observed that, for a trigonometric body force and exact solution, the method suffered considerably from integrating accurately the functions. Only when 10 integration points were set the results were close to the exact solution. 

Another important aspect is that the domain herein used is coincident with the domain in which the Gaussian Quadrature is defined. Otherwise, the method would require the evaluation of the Jacobian to correctly integrate the functions changing the domain of integration. The Boundary Conditions imposed in the problem are also crucial, being the Dirichlet condition null on $\partial \Omega_D$ so that no extra BC methods are implemented. The Neumann BC is simply the inner product of the gradient of the state variable $u$ with the normal vector $\bm{n}$ to the edge of the domain. 

Finally, it is observed that the Least Squares Method indeed yields a more accurate solution than the linearization of the curve fitting the set of data points. Although the method for non-linear systems is more computationally expensive, its employment is justified when the linearization of the curve is not accurate enough.

% ------- BIBLIOGRAPHY -------
\addcontentsline{toc}{section}{References}
\bibliographystyle{abntex2-num}
\bibliography{References}

% ------- APPENDIX -------
\appendix
\section{GitHub Repository}\label{sec:github}
 The source code for this report and every code inhere mentioned can be found in the following GitHub repository: \href{https://github.com/CarlosPuga14/MetodosNumericos_2024S1}{CarlosPuga14/MetodosNumericos\_2024S1}.

\end{document}